\section{Asignación de actividades}

Sub equipos:
\begin{enumerate}
  \item Cortez Vega Jaime Fabian y Lujan De la Rosa Cristian Emir
  \item Arias Mendoza Angel Adrian y Vasquez Moreno Zair De Jesus
  \item Carrasco Pacheco Levi Josue y Nicolas Vasquez Joel Osmar
\end{enumerate}

\vspace{0.5cm}

\begin{longtable}{|P{1.8cm}|P{4.2cm}|P{4.2cm}|P{3.2cm}|}
\hline
\rowcolor{colA}
{\color{white}\textbf{Equipo}} &
{\color{white}\textbf{Actividad}} &
{\color{white}\textbf{Entregas clave}} &
{\color{white}\textbf{Fecha de Entrega}} \\
\hline
\endhead

\rowcolor{colC}
Todos &
Diseño de la estructura general del proyecto, definición de la clase \texttt{Gramatica} y
comprensión del algoritmo de punto fijo. &
 &
Semana 1. \\
\hline

\rowcolor{colD}
1 &
Implementación del algoritmo de PRIMERO y del algoritmo de SIGUIENTE mediante iteración de
punto fijo. &
Módulos \texttt{algoritmo\_primeros.py} y \texttt{algoritmo\_siguientes.py} funcionales con
pruebas sobre gramáticas de ejemplo. &
Semana 2. \\
\hline

\rowcolor{colC}
2 &
Implementación de la clase \texttt{Gramatica} y del lector de archivos \texttt{.txt} en sus
dos formatos (simple y con encabezados). &
Módulos \texttt{gramatica.py} y \texttt{lector\_archivo.py} con manejo de errores y soporte
para $\varepsilon$. &
Semana 2. \\
\hline

\rowcolor{colD}
3 &
Diseño e implementación de la interfaz gráfica: ventana principal, editor de gramática con
numeración de líneas y ventana de resultados con pestañas. &
Módulos de presentación integrados en \texttt{interfaz.py}; interfaz funcional y probada en
Windows y macOS. &
Semana 3. \\
\hline

\end{longtable}

\vspace{0.5cm}

\begin{itemize}
  \item Es importante que cada integrante del equipo registre su avance diario en el informe
        de actividades, describiendo las tareas realizadas o los errores encontrados.
  \item La primera reunión de control de avances tendrá lugar al finalizar la semana 1, donde
        los sub equipos presentarán su diseño y resolverán dudas.
  \item La segunda reunión se realizará al cierre de la semana 2 para integrar los módulos de
        algoritmos y modelo de datos con la interfaz en desarrollo.
  \item Los días restantes se dedicarán a pruebas de integración, corrección de errores y
        presentación final con el profesor.
\end{itemize}
